% MAIN SNAPSHOT — LSPV/RMV::0858::2026-10-02T08:58:00-07:00
% !TEX program = xelatex
\documentclass[11pt]{article}
\usepackage[margin=.86in]{geometry}\usepackage{amsmath,amssymb}
\begin{document}
\section*{LSPV-5 MAIN Snapshot 0858}
\textbf{STATUS: OPEN.} P versus NP remains open. No unrestricted SAT lower bound, arbitrary-algorithm transfer theorem, nonrelativizing hinge, or nonalgebrizing hinge is established.

\section*{CORPUS RECONCILIATION}
\textbf{Corpus surveyed.} Current GitHub history/content index and root \texttt{main.tex}; \texttt{AGENTS.md}, \texttt{CODEX.md}, and the repository pnp-research skill; TASK-001--007; verified MAIN history through 0816; all materially relevant LSPV-1--4 work through MATRIXCERT1, EXCESS1/EXCESSDUAL1, RAMCOMP1, and TRACECLOSE1; HEC/VSC-HEC/G-PPC-HEC/CRS-PPC/CCSPC/RCCN/Search--G; product-span and hostile-control artifacts; PRIME-LQP-1; and indexed GROK/CLAUDE material. Predecessor LSPV-4 has a verified TRACECLOSE1 report, so there is no scientific body to reconstruct. Two stale automation recovery prompts were corrected directly: LSPV-4 now recognizes TRACECLOSE1 as its verified baseline, and LSPV-2 no longer retries the resolved 2026-10-01 missing-write alert.

\textbf{Primary-literature refresh.} Current literature reinforces rather than closes the barriers. Chen--Hu--Ren (ECCC TR25-184) gives new algebrization barriers for circuit lower bounds. Viola (ECCC TR25-059) gives a tight communication lower bound for pointer chasing, but no compiler from unrestricted SAT computations. Jaser--Toran (2025) connect pebbling to restricted algebraic proof systems, again model-specific. Pyne--Tell (ECCC TR25-140, revised 2026) prove unconditional composition time-space tradeoffs for linear-time low-space algorithms but explicitly do not provide the unrestricted polynomial-time SAT transfer needed here. Lin (2026) gives exact Fourier methods for Schur products of constacyclic codes; transfer through the nonlinear standard-binary SU/vMS labeling remains unproved. Raz (ECCC TR26-008) extends Natural-Proofs barriers to explicit linear-function lower bounds under stated cryptographic assumptions, strengthening the guard against treating efficiently recognizable rank properties as a free unrestricted circuit-lower-bound route.

\textbf{New verified results reconciled.}
MATRIXCERT1 gives an exact matrix-native legal-slice certificate
\[
\rho(H)\ge1+\operatorname{rank}X_Q+\sum_i\binom{d_i}{2}.
\]
EXCESS1 corrects the local-nullity narrative:
\[
\dim K_I^{\rm var}=\binom{|I|}{2}-d+\epsilon(I),\qquad
\epsilon(I)=d-\operatorname{rank}G_{E(I)},
\]
so raw local-nullity abundance is retired as a source-specific signal; only excess or finer invariants may distinguish the source. RAMCOMP1 supplies a declared-model reconstruction upper-bound ledger, not a lower bound:
\[
T_E(N)\le N^{c_0K^2+o(1)},\qquad
T_{\rm select}(N)\le N^{\ell+c_0K^2+o(1)}
\]
under its explicit compilation/balance hypotheses. TRACECLOSE1 proves polynomial transcript functionals of a deterministic polynomial-time solver are FP-reacquirable.

BITRANK1 adds a necessary correction to the active product-span route. Complete-orbit exponential-sum methods are not automatically excluded by the short-lag support, but complex/real Gram rank does not imply binary rank. The direct signed-lift/near-orthogonality bridge is therefore rejected. A successful analytic route must end in an intrinsically characteristic-two certificate, such as an odd minor, binary dual-code obstruction, or a model-matched multi-stream correlation-to-rank theorem.

PRIME-LQP-1 formalizes prime-factor provenance. Positional prime coding can injectively encode histories, but ordinary Boolean contexts cannot distinguish extensionally equivalent histories, and deterministic polynomial trace annotations are reacquirable by TRACECLOSE1. Prime magnitude/factorization and direct tag-reading are therefore retired as hardness mechanisms. The only surviving provenance question is whether a naturally required source/witness-bound SAME_REALIZATION relation fails to factor through the semantic quotient and cannot be regenerated or supplied by the maximal semantic/algebraic oracle.

\textbf{Conflicts resolved and stale claims downgraded.}
The prior suggestion that complex Fourier near-orthogonality could directly prove \(\operatorname{rank}_{\mathbb F_2}\Phi\) is downgraded to insufficient. Raw ubiquitous local quadratic relations are no longer evidence of special SU/vMS structure after EXCESS1. Stale missing-write alerts for LSPV-2 and the pre-TRACECLOSE LSPV-4 state are retired. No generic pointer-chasing, pebbling, proof-complexity, streaming, VLSI, or low-space composition theorem is promoted to unrestricted SAT runtime without a model-preserving compiler.

\textbf{Falsifications retained.}
Generic short-source/canonicity/locality/field-degree/list-decoding implications; generic canonical binding; ordinary switching as universal transfer; free atomic shortcuts; optional failed-shortcut costs; raw incidence/tape geometry; causal depth alone; latent branch count; deterministic trace provenance; and real-to-binary rank transfer remain rejected or restricted.

\section*{Shortest surviving proof chain}
The structural branch remains
\[
\boxed{\text{exact source-typed SU/vMS/Hadamard family}
\to \text{binary structural growth}
\to \text{charged SAME\_REALIZATION acquisition}
\to \text{model-matched reconstruction}
\to \text{nonrelativizing/nonalgebrizing contradiction}.}
\]
The causal branch remains a hostile-control laboratory:
\[
\text{HEC/CRS-PPC/CCSPC/RCCN/Search--G}
\to \text{restricted no-reread/state theorems},
\]
but static reread, rerun, memoization, transcript closure, and polynomial selector reacquisition block generic promotion.

\section*{Barrier ledger}
\textbf{OPEN:} exact reconstruction-compatible source typing; asymptotic binary product-span/exterior growth on that exact family; arbitrary-algorithm acquisition/necessity; exact quantitative reconstruction parameters; and a genuine nonrelativizing plus nonalgebrizing terminal step. Any one OPEN arrow keeps P versus NP OPEN.

\section*{Adopted / rejected handoffs}
\textbf{ADOPT:} MATRIXCERT1 certificate pipeline; EXCESS1 baseline subtraction; RAMCOMP1 only as conservative declared-model accounting; BITRANK1 odd-minor/binary-dual reformulation; PRIME-LQP only as a hostile SAME_REALIZATION experiment.

\textbf{REJECT/RETIRE:} complex Gram rank as binary rank; raw local-nullity abundance; prime factorization or deterministic prime trace tags as hardness; generic communication/pebbling/proof-system transfers without a compiler.

\section*{Single highest-value team target}
For one exact reconstruction-compatible SU/vMS polynomial-bit orbit, construct the source-transcribed binary matrix \(\Phi\) and legal slice matrices \(X_s\), then prove or falsify
\[
\boxed{\operatorname{rank}_{\mathbb F_2}\Phi=\omega(m)}
\]
using an intrinsically characteristic-two certificate. Compute in the same pipeline \(\operatorname{rank}_{\mathbb Q}\Phi\), odd/maximal minors, binary dual annihilators, MATRIXCERT1 pair incidence, and EXCESS1 profiles, with COVERBREAK1 as the mandatory hostile control. Do not advance to provenance as a proof hinge unless this structural arrow survives and the downstream SAME_REALIZATION relation is independently necessary.

\section*{What this does not prove}
No finite experiment, coding-theory identity, provenance encoding, restricted causal lower bound, reconstruction upper-bound ledger, or current literature result proves an unrestricted SAT lower bound. P versus NP remains OPEN.

\begin{thebibliography}{9}
\bibitem{CHR25} L. Chen, Y. Hu, H. Ren, New Algebrization Barriers to Circuit Lower Bounds via Communication Complexity of Missing-String, ECCC TR25-184, 2025.
\bibitem{Viola25} E. Viola, Communication Complexity of Pointer Chasing via the Fixed-Set Lemma, ECCC TR25-059, 2025.
\bibitem{JT25} L.-M. Jaser, J. Toran, Pebble Games and Algebraic Proof Systems, 2025.
\bibitem{PT26} E. Pyne, R. Tell, Composing Low-Space Algorithms, ECCC TR25-140, rev. 2026.
\bibitem{Lin26} P. Lin, Schur Products of Constacyclic Codes via the Constacyclic Discrete Fourier Transform, 2026.
\bibitem{Raz26} R. Raz, A Note on Natural-Proofs for Super-Linear Lower Bounds for Linear Functions, ECCC TR26-008, rev. 2026.
\end{thebibliography}
\end{document}
% END_OF_VERIFIED_MAIN::LSPV-5::0858
